%! AP Statistics — Unit 2, Lesson 2.1 — Group Activity
\documentclass[10pt]{article}
\usepackage{apstats-article}
\usepackage{apstats-boxes}

%=============================================================================
\begin{document}

\pageheader{Unit 2, Lesson 2.1}{Group Activity}
\namepartnerperiod

\vspace{0.1in}

% ── Instructions ─────────────────────────────────────────────────────────────
\begin{tcolorbox}[colback=sky, colframe=navy, boxrule=0.8pt, arc=2mm,
    left=3mm, right=3mm, top=2mm, bottom=2mm]
\small
For each study below: (1) identify the two variables and their types,
(2) decide whether the apparent association is likely \textbf{real} or possibly \textbf{random},
(3) name one possible \textbf{confounding variable}, and (4) state whether the study
could establish \textbf{causation}.
\end{tcolorbox}

\vspace{0.14in}

% ════════════════════════════════════════════════════════════════════════════
% SCENARIO 1
% ════════════════════════════════════════════════════════════════════════════
\begin{scenariobox}[Scenario 1]{navylight}

\begin{headlinebox}{sky}
\small\textit{``A study of 1{,}500 adults found that those who own pets report higher
levels of happiness than those who do not own pets.''}
\end{headlinebox}

\vspace{8pt}
\renewcommand{\arraystretch}{1.8}
\begin{tabularx}{\linewidth}{@{} l X @{}}
\textbf{Two variables:}          & \blank{9cm} \\
\textbf{Variable types:}         & \blank{9cm} \\
\textbf{Real or possibly random?}  & \blank{9cm} \\
\textbf{One confounding variable:} & \blank{9cm} \\
\textbf{Establishes causation?}  & \blank{9cm} \\
\end{tabularx}

\end{scenariobox}

\vspace{0.16in}

% ════════════════════════════════════════════════════════════════════════════
% SCENARIO 2
% ════════════════════════════════════════════════════════════════════════════
\begin{scenariobox}[Scenario 2]{greenacc}

\begin{headlinebox}{greenbg}
\small\textit{``Researchers found that states with higher ice cream sales also have
higher drowning rates.''}
\end{headlinebox}

\vspace{8pt}
\renewcommand{\arraystretch}{1.8}
\begin{tabularx}{\linewidth}{@{} l X @{}}
\textbf{Two variables:}          & \blank{9cm} \\
\textbf{Variable types:}         & \blank{9cm} \\
\textbf{Real or possibly random?}  & \blank{9cm} \\
\textbf{One confounding variable:} & \blank{9cm} \\
\textbf{Establishes causation?}  & \blank{9cm} \\
\end{tabularx}

\vspace{8pt}
\small\textbf{Why is this association real but \emph{not} causal? What is the confounding variable?}\\[4pt]
\writeline\\[1pt]\writeline

\end{scenariobox}

\vspace{0.16in}

% ════════════════════════════════════════════════════════════════════════════
% SCENARIO 3
% ════════════════════════════════════════════════════════════════════════════
\begin{scenariobox}[Scenario 3 \normalfont\small\color{white!80!black}\quad stretch]{redacc}

\begin{headlinebox}{redbg}
\small\textit{``A study of 12 high school students found that the 4 students who slept
the most also had the highest grades. The 4 students who slept the least had the lowest grades.''}
\end{headlinebox}

\vspace{8pt}
\renewcommand{\arraystretch}{1.8}
\begin{tabularx}{\linewidth}{@{} l X @{}}
\textbf{Two variables:}          & \blank{9cm} \\
\textbf{Variable types:}         & \blank{9cm} \\
\textbf{Real or possibly random?}  & \blank{9cm} \\
\textbf{One confounding variable:} & \blank{9cm} \\
\textbf{Establishes causation?}  & \blank{9cm} \\
\end{tabularx}

\vspace{8pt}
\small\textbf{Why does the small sample size ($n = 12$) make this association hard to trust?}\\[4pt]
\writeline\\[1pt]\writeline\\[1pt]\writeline

\end{scenariobox}

\vspace{0.16in}

% ── Reflection ───────────────────────────────────────────────────────────────
\begin{tcolorbox}[colback=goldbg, colframe=goldacc, boxrule=0.8pt, arc=2mm,
    left=3mm, right=3mm, top=2mm, bottom=2mm]
\small\textbf{Reflection \textnormal{\color{slate}(one response per pair):}}
\quad Which scenario best illustrates why ``association $\neq$ causation''?
Explain in one or two sentences.\\[6pt]
\writeline\\[1pt]\writeline
\end{tcolorbox}

\end{document}
